\chapter{Inventory Models}\label{ch:hf:inventory-models}

A market maker who has just bought three times in a row owns stock nobody has asked for yet. Lowering both quotes a little makes the next trade more
likely to be a sale; lowering them too much gives the spread away. Avellaneda and Stoikov turned the trade-off into a control problem whose
approximate solution fits on one line, and simulated it: at their benchmark risk aversion the inventory-aware quotes earned 5\% less than symmetric
quotes of the same average width, with half the standard deviation. This chapter solves the problem, reproduces their experiment, solves the variant
with a running \omterm{def:hf:inventory-models:penalty}{inventory penalty} exactly, and then takes the answer to a large-tick market, where a quote can only be at the touch or absent: there,
on this chapter's simulated sessions, controlling inventory turns a losing quoter into a break-even one.

\section{The market maker's control problem}

Ho and Stoll set the problem in 1981: a dealer whose prices depend on its inventory because holding inventory is risky. The modern form keeps the
ingredients and makes the order flow explicit. The mid is $S_t=S_0+\sigma W_t$ (arithmetic, over a trading session). The market maker quotes a bid
$S^b_t=S_t-\delta^b_t$ and an ask $S^a_t=S_t+\delta^a_t$, where the quote depths $\delta^b_t,\delta^a_t$ are its controls (local symbols of this
chapter; $\delta$ alone stays the accrual fraction). Market sell orders reach the bid, and market buys the ask, at rates that fall with the depth.
Its inventory $q_t$ rises by one on each bid fill and falls by one on each ask fill; its cash $X_t$ moves by the fill prices.

\begin{definition}[Fill intensity]\label{def:hf:inventory-models:intensity}
The \emph{fill intensity}\index{fill intensity} of a resting quote is the rate at which it is executed, as a function of its depth $\delta$ from the
mid. The inventory models assume $\lambda(\delta)=A\,e^{-k\delta}$, the same on both sides, with $A$ and $k$ constant (local symbols: $A$ is not an
annuity and $k$ not a log-moneyness here).
\end{definition}

The exponential form is not a law. It comes from assuming that market orders arrive at a constant rate with sizes whose price impact reaches a
depth $\delta$ with probability $e^{-k\delta}$; it is convenient because it makes the control problem solvable, and it is testable, since a quoter can
measure how often it is filled at each depth (\cref{sec:hf:inventory-models:limits}).

\section{The Avellaneda--Stoikov model}

\begin{definition}[Avellaneda--Stoikov model]\label{def:hf:inventory-models:as}
The \emph{Avellaneda--Stoikov model}\index{Avellaneda--Stoikov model} is the market maker's problem with an arithmetic Brownian mid, exponential fill
intensities, a horizon $T$ and an exponential utility of terminal wealth: the market maker chooses $\delta^b_t,\delta^a_t$ to maximise
$\E\bigl[-\exp\bigl(-\gamma\,(X_T+q_TS_T)\bigr)\bigr]$, with $\gamma>0$ its risk aversion.
\end{definition}

The Hamilton--Jacobi--Bellman equation of One Quant Book 4, chapter 9 applies. With the ansatz $u(t,x,s,q)=-\exp(-\gamma(x+\theta(t,s,q)))$, the
first-order conditions give the optimal depths as functions of $\theta$'s differences across inventories; Avellaneda and Stoikov then expand
$\theta$ to second order in $q$ and obtain the approximate solution below.

\section{Reservation price and optimal spread}

\begin{definition}[Reservation price]\label{def:hf:inventory-models:reservation}
The \emph{reservation price}\index{reservation price} of a market maker holding $q$ is the price at which it is indifferent to its current position
and one more or one fewer unit; in the \omterm{def:hf:inventory-models:as}{Avellaneda--Stoikov model}, approximately
\[
S^{\mathrm{res}}_t=S_t-q_t\,\gamma\sigma^2\,(T-t).
\]
\end{definition}

\begin{proposition}[The Avellaneda--Stoikov quotes]\label{prop:hf:inventory-models:as}
To second order in inventory, the optimal quotes are centred on the \omterm{def:hf:inventory-models:reservation}{reservation price}, $S^{b,a}_t=S^{\mathrm{res}}_t\mp\tfrac12\Delta_t$, with total
width
\[
\Delta_t=\delta^a_t+\delta^b_t=\gamma\sigma^2\,(T-t)+\frac{2}{\gamma}\ln\Bigl(1+\frac{\gamma}{k}\Bigr).
\]
\end{proposition}

\admitted
The derivation is in Avellaneda and Stoikov's paper; \cref{prop:hf:inventory-models:cj} below derives an exact solution of a close relative.

Two effects are separated. The \omterm{def:hf:inventory-models:reservation}{reservation price} moves against the inventory in proportion to risk aversion, variance and the time left: a long
market maker quotes lower on both sides, selling more readily and buying less. The width has an inventory-risk part, $\gamma\sigma^2(T-t)$, and a
monopoly part, $(2/\gamma)\ln(1+\gamma/k)$, which tends to $2/k$ as $\gamma\to0$: with no risk aversion the market maker still quotes a spread,
set by how fast fills fall off with depth. The skew of the quotes is the quote skewing of One Quant Book 2, chapter 15, derived from a model instead
of a rule of thumb.

\begin{definition}[Inventory penalty]\label{def:hf:inventory-models:penalty}
An \emph{inventory penalty}\index{inventory penalty} replaces risk aversion in utility by a cost charged on the position while it is held: the
market maker maximises $\E\bigl[X_T+q_TS_T-\alpha q_T^2-\phi\int_0^Tq_t^2\,dt\bigr]$, with a running penalty $\phi$ and a terminal one $\alpha$.
\end{definition}

The running penalty is a risk charge the desk chooses, not a preference it has. With $\phi=\tfrac12\gamma\sigma^2$ it charges exactly the variance
that the position adds per unit of time; a desk can also set it from a value-at-risk budget.

\begin{proposition}[Exact solution with an inventory penalty]\label{prop:hf:inventory-models:cj}
With inventory confined to $\{-\bar q,\dots,\bar q\}$, write the value function as $x+qs+h(t,q)$ and put $h=\tfrac1k\ln\omega$. Then
$\omega(t)=\exp\bigl(M(T-t)\bigr)z$, with $z_q=e^{-\alpha kq^2}$ and $M$ tridiagonal: $M_{q,q}=-\phi kq^2$ and $Ae^{-1}$ next to the diagonal.
The optimal depths are
\[
\delta^a(t,q)=\frac1k+h(t,q)-h(t,q-1),\qquad \delta^b(t,q)=\frac1k+h(t,q)-h(t,q+1),
\]
and the side that would breach the bound is not quoted.
\end{proposition}

\begin{proof}
The Hamilton--Jacobi--Bellman equation reduces to $\partial_th-\phi q^2+A\sup_{\delta}e^{-k\delta}\bigl(\delta+h(t,q-1)-h(t,q)\bigr)+A\sup_\delta
e^{-k\delta}\bigl(\delta+h(t,q+1)-h(t,q)\bigr)=0$, with $h(T,q)=-\alpha q^2$; the mid's diffusion drops out because the value is linear in $s$. The
first-order condition of the first supremum gives $\delta^a=1/k+h(t,q)-h(t,q-1)$ and its value $\tfrac1ke^{-1}e^{k(h(t,q-1)-h(t,q))}$, and likewise
on the bid. Substituting $h=\tfrac1k\ln\omega$ turns the equation into $\partial_t\omega_q=\phi kq^2\omega_q-Ae^{-1}(\omega_{q-1}+\omega_{q+1})$, the
linear system $\partial_t\omega+M\omega=0$ with $\omega(T)=z$.
\end{proof}

\Cref{fig:hf:inventory-models:depths} shows the solution for the chapter's parameters. Flat, the market maker quotes symmetric depths above $1/k$;
long, it lowers the ask depth and raises the bid depth, and at large inventories the ask depth becomes negative: the model tells the market maker to
offer below the mid, which in a real book means crossing the spread. The model has no notion of a spread to cross; chapter 4 bounds the inventory
instead.

\begin{omfigure}
\begin{tikzpicture}
\begin{axis}[width=11cm,height=5.4cm,xlabel={\small inventory $q$},ylabel={\small depth from the mid},tick label style={font=\footnotesize},
  xmin=-10,xmax=10,ymin=-3,ymax=4.5,grid=major,grid style={gray!25},table/col sep=comma,
  legend style={font=\scriptsize,at={(0.5,-0.3)},anchor=north},legend columns=3]
\addplot[omDef,thick] table[x=q,y=b05]{figdata/hft/03-inventory-models/depths.csv};
\addplot[omDef,thick,dashed] table[x=q,y=a05]{figdata/hft/03-inventory-models/depths.csv};
\addplot[omThm,thick] table[x=q,y=b5]{figdata/hft/03-inventory-models/depths.csv};
\addplot[omThm,thick,dashed] table[x=q,y=a5]{figdata/hft/03-inventory-models/depths.csv};
\addplot[black,thick] table[x=q,y=b50]{figdata/hft/03-inventory-models/depths.csv};
\addplot[black,thick,dashed] table[x=q,y=a50]{figdata/hft/03-inventory-models/depths.csv};
\draw[gray] (axis cs:-10,0) -- (axis cs:10,0);
\legend{{bid, $\phi=0.5$},{ask, $\phi=0.5$},{bid, $\phi=5$},{ask, $\phi=5$},{bid, $\phi=50$},{ask, $\phi=50$}}
\end{axis}
\end{tikzpicture}

\omcaption{Optimal bid depth (solid) and ask depth (dashed) at the start of the session against inventory, for three running penalties $\phi$
(\cref{prop:hf:inventory-models:cj}; $A=140$, $k=1.5$, $T=1$, $\alpha=1$, inventory within $\pm30$). Below zero the ask is offered under the mid.
Data: \texttt{firm.invmm.CJSolution}.}\label{fig:hf:inventory-models:depths}
\end{omfigure}

\section{The Avellaneda--Stoikov experiment, reproduced}

Avellaneda and Stoikov simulated their quotes with $S_0=100$, $T=1$, $\sigma=2$, time steps of 0.005, $k=1.5$ and $A=140$, the mid moving up or down
by $\sigma\sqrt{0.005}$ at each step, and compared them with a symmetric strategy that quotes the same average width around the mid. The same
experiment, run with \texttt{firm.invmm} on 1,000 paths with the two strategies on common random numbers, gives:

\begin{center}
\small
\begin{tabular}{@{}llcccccc@{}}
\toprule
$\gamma$ & strategy & width & profit & sd (profit) & final $q$ & sd (final $q$) & paper: profit, sd\\
\midrule
0.01 & inventory & 1.35 & 68.45 & 9.06 & $-0.05$ & 5.24 & 68.6, 8.7\\
 & symmetric & 1.35 & 68.50 & 13.91 & $-0.05$ & 8.99 & 68.8, 12.8\\
0.1 & inventory & 1.49 & 64.83 & 6.49 & 0.09 & 2.97 & 65.0, 6.6\\
 & symmetric & 1.49 & 68.15 & 14.00 & $-0.21$ & 8.66 & 68.4, 12.7\\
1 & inventory & 3.03 & 31.44 & 4.83 & 0.05 & 1.66 & 31.4, 5.0\\
 & symmetric & 3.03 & 43.30 & 10.34 & $-0.14$ & 5.39 & 44.0, 11.0\\
\bottomrule
\end{tabular}
\end{center}

The reproduction agrees with the published tables to within sampling error. Near risk neutrality the two strategies earn the same and the
inventory strategy is already less variable; at $\gamma=0.1$ it gives up 5\% of the profit for half the standard deviation; at $\gamma=1$ it gives
up 27\%. \Cref{fig:hf:inventory-models:path} redraws one path: the quotes follow the \omterm{def:hf:inventory-models:reservation}{reservation price}, which leaves the mid whenever inventory
builds up.

\begin{omfigure}
\begin{tikzpicture}
\begin{axis}[width=11cm,height=5.4cm,xlabel={\small time},ylabel={\small price},tick label style={font=\footnotesize},xmin=0,xmax=1,
  grid=major,grid style={gray!25},table/col sep=comma,legend style={font=\scriptsize,at={(0.5,-0.3)},anchor=north},legend columns=4]
\addplot[black,thick] table[x=t,y=s]{figdata/hft/03-inventory-models/path.csv};
\addplot[omProp,thick,dashed] table[x=t,y=r]{figdata/hft/03-inventory-models/path.csv};
\addplot[omDef,thin] table[x=t,y=bid]{figdata/hft/03-inventory-models/path.csv};
\addplot[omThm,thin] table[x=t,y=ask]{figdata/hft/03-inventory-models/path.csv};
\legend{mid,reservation price,bid,ask}
\end{axis}
\end{tikzpicture}

\omcaption{One path of the Avellaneda--Stoikov experiment ($\gamma=0.1$): the mid, the \omterm{def:hf:inventory-models:reservation}{reservation price} and the inventory strategy's quotes. The
width narrows as the horizon approaches. Data: \texttt{hf\_inventory.one\_path}.}\label{fig:hf:inventory-models:path}
\end{omfigure}

With a running penalty instead of utility, the exact quotes of \cref{prop:hf:inventory-models:cj} trace a frontier of mean against standard
deviation of profit that lies well above the symmetric quotes' (\cref{fig:hf:inventory-models:frontier}): a small penalty, $\phi=0.5$, keeps
97\% of the profit of the best symmetric width (66.7 against 68.5) with half its standard deviation (6.9 against 13.9).

\begin{omfigure}
\begin{tikzpicture}
\begin{axis}[width=11cm,height=5.4cm,xlabel={\small standard deviation of profit},ylabel={\small mean profit},tick label style={font=\footnotesize},
  xmin=4,xmax=15,ymin=25,ymax=72,grid=major,grid style={gray!25},table/col sep=comma,
  legend style={font=\scriptsize,at={(0.5,-0.3)},anchor=north},legend columns=3]
\addplot[omThm,thick,mark=*,mark size=1.3pt] table[x=sd,y=profit]{figdata/hft/03-inventory-models/frontier_cj.csv};
\addplot[omDef,thick,dashed,mark=square*,mark size=1.3pt] table[x=sd,y=profit]{figdata/hft/03-inventory-models/frontier_sym.csv};
\addplot[black,only marks,mark=triangle*,mark size=2pt] table[x=sd,y=profit]{figdata/hft/03-inventory-models/frontier_as.csv};
\legend{{penalty $\phi$ (exact)},{symmetric width},{Avellaneda--Stoikov $\gamma$}}
\end{axis}
\end{tikzpicture}

\omcaption{Mean against standard deviation of profit over 1,000 paths of the Avellaneda--Stoikov experiment: the exact quotes with running penalties
$\phi$ from 0 to 50, symmetric quotes with half-widths from 0.5 to 2, and the three risk aversions of the table. Data:
\texttt{hf\_inventory.frontier}.}\label{fig:hf:inventory-models:frontier}
\end{omfigure}

\section{The limits of the closed forms}\label{sec:hf:inventory-models:limits}

Four assumptions carry the model, and each fails somewhere.

\begin{remark}[What the model leaves out]\label{rem:hf:inventory-models:limits}
(i) \emph{Continuous prices.} Quotes live on a tick grid; the model's depths are real numbers. (ii) \emph{Exponential intensities.} Fills depend
on queue position and on the book behind the quote, not on depth alone. (iii) \emph{No information.} The mid is a martingale unaffected by the fills:
there is no adverse selection, so a fill is never bad news (chapter 4 adds it). (iv) \emph{Constant parameters.} Volatility and order flow vary
through the day and with news.
\end{remark}

The first two failures are severe in a large-tick market. On Book 7's simulated stock, an exploring quoter that moved its one-lot quotes at random
among the touch and one, two or three ticks behind it, every five seconds for an hour, spent about 1,640 seconds at the touch and got 8 fills, and
about 5,280 seconds behind it and got none: the book's depth behind the touch is never reached by the market orders of this market. The only
choice is to rest at the touch or not.

The problem remains a control problem. With $\lambda$ the \omterm{def:hf:inventory-models:intensity}{fill intensity} of one lot at the touch per side and $c$ the capture per fill net of
adverse selection, the value $h(t,q)$ solves
\[
\partial_th-\phi q^2+\lambda\max\bigl(0,c+h(t,q-1)-h(t,q)\bigr)+\lambda\max\bigl(0,c+h(t,q+1)-h(t,q)\bigr)=0,
\]
and the market maker rests on a side exactly when that side's bracket is positive (\cref{lst:hf:inventory-models:touch}). The quote never moves
away from the touch; the skew of the continuous model becomes a decision about which sides to show.

\section{Tutorial: from the model to a quoter}\label{tut:hf:inventory-models:tut}

\begin{tutorial}
\textbf{Goal.} Solve the model, reproduce the published experiment, then run the touch-only version on the simulated market. \textbf{End state:} the
table above, \cref{fig:hf:inventory-models:depths,fig:hf:inventory-models:path,fig:hf:inventory-models:frontier}, and the table below.
\begin{tutsteps}
\item \textbf{Solve.} \texttt{CJSolution} builds $M$, steps $\omega$ back from $T$ with one matrix exponential, renormalising at every step (a
constant factor drops out of every depth, and without it $e^{-\alpha kq^2}$ underflows at $q=30$).
\omcode{code/firm/invmm/firm_invmm.py}{55}{75}{The exact solution with a running penalty (\cref{prop:hf:inventory-models:cj}).}\label{lst:hf:inventory-models:cj}
\item \textbf{Reproduce.} \texttt{hf\_inventory.as\_table()} runs the inventory and symmetric strategies on 1,000 common paths for each $\gamma$.
\item \textbf{Calibrate for the tape.} On three calibration sessions of twenty minutes, the chapter 1 symmetric Quoter, always at the touch, was
filled 241 times: $\lambda=0.033$ lots a second per side, and its fills marked out, ten seconds later, at $c=0.139$ ticks a share.
\item \textbf{Solve the touch-only problem} backwards on a one-second grid and quote it (\cref{lst:hf:inventory-models:touch}).
\omcode{code/firm/invmm/firm_invmm.py}{89}{110}{The touch-only (large-tick) market maker: rest on a side only if its bracket is positive.}\label{lst:hf:inventory-models:touch}
\item \textbf{Compare} on six twenty-minute sessions of the simulated market with \texttt{tape\_compare(phi)}.
\end{tutsteps}
\textbf{What to change next.} Let $c$ depend on the book's imbalance (chapter 7); estimate $\lambda$ by time of day.
\end{tutorial}

\begin{center}
\small
\begin{tabular}{@{}lccccc@{}}
\toprule
quoter (one lot, six sessions) & P\&L (\$) & sd (\$) & shares & mean $|q|$ (lots) & mean max $|q|$\\
\midrule
symmetric, at most 5 lots & $-65.17$ & 97.87 & 15,783 & 2.51 & 5.00\\
touch-only, $\phi=3\times10^{-4}$ & $-0.75$ & 29.53 & 14,317 & 1.12 & 3.67\\
touch-only, $\phi=10^{-3}$ & 11.00 & 8.41 & 11,433 & 0.57 & 2.33\\
\bottomrule
\end{tabular}
\end{center}

The penalised quoters trade almost as much and hold a fraction of the inventory; they lose less, and vary far less. With six sessions the means are
imprecise (the symmetric quoter's standard error is \$40), and the direction is the lesson rather than the size: in this market a filled quote is
followed on average by an adverse move, so the inventory the symmetric quoter accumulates is where it loses.

\section{Build: the inventory-model library}\label{bld:hf:inventory-models:invmm}

\begin{build}
\bfield{Purpose} Inventory-aware quotes from closed forms and exact solutions, and a fast simulator to compare quoting policies on common random
numbers.
\bfield{Interface} \texttt{as\_quotes(s, q, tau, gamma, sigma, k)}, \texttt{as\_policy}; \texttt{CJSolution(A, k, phi, alpha, qmax, T, n).depths(t, q)},
\texttt{cj\_policy}; \texttt{TouchSolution(lam, c, phi, alpha, qmax, T, n).post(t, q)}; \texttt{estimate\_intensity(depth, time, fills)};
\texttt{simulate(policy, T, dt, sigma, A, k, paths, seed, qmax, mid)}; \texttt{symmetric(half)}.
\bfield{Rules} Policies see only time and inventory. The simulator's random numbers depend on the seed and the step only, so policies are compared on
the same paths. A side at the inventory bound is never quoted.
\bfield{Acceptance tests} \texttt{code/firm/invmm/tests/}: the closed form by hand; symmetric depths when flat and mirror-image skews; the monopoly depth
$1/k$ far from the horizon without penalties; the \omterm{def:hf:inventory-models:intensity}{fill intensity} recovered from synthetic counts; the simulator deterministic and its fill count
equal to $2Ae^{-k\delta}T$; the touch-only solution posting less when loaded and never beyond the bound.
\bfield{Stretch} A two-dimensional grid for two correlated instruments (chapter 4); intensities depending on queue position (chapter 5).
\end{build}

\begin{omsources}
\item M.~Avellaneda and S.~Stoikov, ``High-frequency trading in a limit order book'', \emph{Quantitative Finance} 8(3), 2008.
\item T.~Ho and H.~R.~Stoll, ``Optimal dealer pricing under transactions and return uncertainty'', \emph{Journal of Financial Economics} 9(1), 1981.
\item \'A.~Cartea, S.~Jaimungal and J.~Penalva, \emph{Algorithmic and High-Frequency Trading}, Cambridge University Press, 2015.
\end{omsources}

\section{Exercises}

\begin{exercise}[$\star$]\label{exo:hf:inventory-models:1}
A market maker with $\gamma=0.1$, $\sigma=2$ and half the session left ($T-t=0.5$) holds 3 units at a mid of 100. Compute its \omterm{def:hf:inventory-models:reservation}{reservation price}.
\end{exercise}

\begin{exercise}[$\star$]\label{exo:hf:inventory-models:2}
With $k=1.5$, compute the Avellaneda--Stoikov width at the start of the session ($T-t=1$) for $\gamma=0.1$, and its monopoly part.
\end{exercise}

\begin{exercise}[$\star$]\label{exo:hf:inventory-models:3}
A quote rests at depth 0.5 for 1,000 seconds and is filled 40 times; at depth 1.0 for 1,000 seconds it is filled 18 times. Estimate $A$ and $k$.
\end{exercise}

\begin{exercise}[$\star\star$]\label{exo:hf:inventory-models:4}
Show that the Avellaneda--Stoikov width tends to $2/k$ as $\gamma\to0$ at the horizon, and explain why a risk-neutral market maker still quotes a
spread.
\end{exercise}

\begin{exercise}[$\star\star$]\label{exo:hf:inventory-models:5}
From \cref{fig:hf:inventory-models:depths}, at which inventory does the ask depth turn negative for $\phi=5$, and what would the quote mean in a
real book?
\end{exercise}

\begin{exercise}[$\star\star$]\label{exo:hf:inventory-models:6}
Why does the symmetric strategy earn more on average than the inventory strategy at $\gamma=0.1$, and why is that no reason to prefer it?
\end{exercise}

\begin{exercise}[$\star\star\star$]\label{exo:hf:inventory-models:7}
\emph{Coding.} With \texttt{CJSolution(140, 1.5, 5.0, 1.0, 30, 1.0, 200)}, give the depths at the start of the session for $q=0$ and for $q=3$.
\end{exercise}

\begin{exercise}[$\star\star\star$]\label{exo:hf:inventory-models:8}
\emph{Find the flaw.} ``We fitted $A$ and $k$ on our quotes one and two ticks behind the touch in a one-tick stock and got $k=0$: fills do not depend
on depth, so we will quote deeper and earn more.''
\end{exercise}

\section{Problem: Three Buys in a Row}

\begin{problem}[{Weekend problem --- three buys in a row}]\label{pb:hf:inventory-models:1}
A market maker has just bought three times. Decide what to quote next, first in the textbook model, then in the market it actually trades.

\medskip\noindent\textbf{Part I --- The model.}
\begin{enumerate}
\item State the \omterm{def:hf:inventory-models:as}{Avellaneda--Stoikov model}: dynamics, controls and objective.
\item Define the \omterm{def:hf:inventory-models:intensity}{fill intensity} and give the model's assumption.
\item Define the \omterm{def:hf:inventory-models:reservation}{reservation price} and write its approximate form.
\item Write the approximate width and split it into its two parts.
\item Compute the \omterm{def:hf:inventory-models:reservation}{reservation price} and quotes at $q=3$, $\gamma=0.1$, $\sigma=2$, $k=1.5$, $T-t=0.5$, $S=100$.
\end{enumerate}

\medskip\noindent\textbf{Part II --- The exact solution.}
\begin{enumerate}[resume]
\item Define an \omterm{def:hf:inventory-models:penalty}{inventory penalty} and relate the running penalty to risk aversion.
\item Derive the depths of \cref{prop:hf:inventory-models:cj} from the Hamilton--Jacobi--Bellman equation.
\item Why does the mid's volatility drop out of the equation for $h$?
\item Give the depths at $q=0$ and $q=3$ for $\phi=5$, and interpret the negative ask depth at large inventory.
\end{enumerate}

\medskip\noindent\textbf{Part III --- The experiment.}
\begin{enumerate}[resume]
\item Give the experiment's parameters.
\item Give profit and standard deviation of both strategies at the three risk aversions, and compare with the paper.
\item At which risk aversion do the two strategies' profits diverge, and by how much?
\item What does the penalty frontier add?
\end{enumerate}

\medskip\noindent\textbf{Part IV --- The market it trades.}
\begin{enumerate}[resume]
\item What did the exploring quoter find behind the touch, and why?
\item Write the touch-only equation and its decision rule.
\item How were $\lambda$ and $c$ calibrated, and why on separate sessions?
\item State the \emph{named result}: the P\&L and inventory of the symmetric and the penalised quoters on the simulated market, and, in the model, the
profit and standard deviation of the two strategies at $\gamma=0.1$.
\item Why are six sessions enough for the direction and not for the size?
\item Which assumption of the model does the simulated market violate most?
\item In one sentence: what should a market maker do after three buys in a row?
\end{enumerate}
\end{problem}

\section{Interview questions}

\begin{interviewq}[$\star$ \iqroles{trader}]\label{iq:hf:inventory-models:1}
You are long and want to get flat. Do you move your bid, your ask, or both, and which way?
\end{interviewq}

\begin{interviewq}[$\star\star$ \iqroles{researcher}]\label{iq:hf:inventory-models:2}
Derive the optimal depth of a risk-neutral market maker facing a \omterm{def:hf:inventory-models:intensity}{fill intensity} $Ae^{-k\delta}$, ignoring inventory.
\end{interviewq}

\begin{interviewq}[$\star\star$ \iqroles{researcher}]\label{iq:hf:inventory-models:3}
Why does the Avellaneda--Stoikov \omterm{def:hf:inventory-models:reservation}{reservation price} scale with $T-t$, and is that sensible for a market maker that trades every day?
\end{interviewq}

\begin{interviewq}[$\star\star$ \iqroles{developer}]\label{iq:hf:inventory-models:4}
Your matrix-exponential solver returns NaN depths at large inventories. What happened, and how do you fix it without changing the answer?
\end{interviewq}

\begin{interviewq}[$\star\star$ \iqroles{risk}]\label{iq:hf:inventory-models:5}
How would you choose the running \omterm{def:hf:inventory-models:penalty}{inventory penalty} for a desk with a daily loss limit?
\end{interviewq}

\begin{interviewq}[$\star\star\star$ \iqroles{researcher}]\label{iq:hf:inventory-models:6}
In a one-tick stock, what replaces the choice of depth, and how would you set up the control problem?
\end{interviewq}
